\documentclass[conference]{IEEEtran}

\usepackage[english]{babel}
\usepackage{amsmath}
\usepackage{graphicx}
\usepackage{url}
\usepackage{hyperref}
\usepackage{microtype}
\hypersetup{hidelinks}

\title{Janus: Evidence-Based Governance for AI-Assisted Development Systems\\
\large Version 1.1 Candidate}

\author{%
  \IEEEauthorblockN{Martín Nicolás Sánchez Morales}\\
  \IEEEauthorblockA{Janus Governance Project\\
  Argentina}%
}

\begin{document}

\maketitle

\begin{abstract}
AI-assisted development introduces structural challenges in traceability, authority,
and reproducibility. This paper presents Janus, a protocol-driven governance layer
that records and signals deviations from evidence-based workflow protocols across
development environments.

Rather than measuring system performance, this work evaluates governance through
conformance between protocol-defined expectations and runtime-observed behavior.

We specify an evidence-based governance benchmark and delimit the evidence
required before empirical results may be reported.
\end{abstract}

\begin{IEEEkeywords}
Governance, Observability, AI-assisted development, Traceability, Runtime systems
\end{IEEEkeywords}


\section{Introduction}


This work extends the Janus foundational model introduced in [1].

AI-assisted systems introduce opacity into decision-making processes, making it
difficult to trace actions, verify authority, and reconstruct execution history.

Janus introduces a governance layer that records and signals deviations from
protocol-defined workflows, providing:
\begin{itemize}
    \item Explicit evidence (E\textsuperscript{+} and E\textsuperscript{-})
    \item Human authority boundaries
    \item Deterministic reconstruction of governance state
\end{itemize}

Rather than modifying underlying systems, Janus wraps workflows with a protocol
stack that monitors traceability and signals accountability deviations at the operational level.

\section{Architecture}

Janus operates as a protocol-driven governance layer composed of four distinct
components:

\begin{itemize}
    \item \textbf{Core Layer:} defines events, evidence types, and authority boundaries
    \item \textbf{Protocol Layer:} defines operational governance rules
    \item \textbf{Runtime Layer:} captures telemetry and execution events
    \item \textbf{Workflow Layer:} mediates interaction between agents and human authority
\end{itemize}

The Protocol, Runtime, and Workflow layers are subordinate to Janus Core Lite
v0.2. Domain activity maps to \texttt{MANAGEMENT\_LOG}; schema evolution maps to
\texttt{SCHEMA\_LOG}; and governance events, accountability records, and audit
findings map to \texttt{AUDIT\_LOG}. Canonical governance events are appended to
\texttt{AUDIT\_LOG} exclusively through the Audit Writer.

\section{Governance Invariants}

Janus defines four operational invariants that govern system behavior:

\medskip
\noindent\textbf{Invariant 1 (Evidence Completeness):}
For every governance-relevant expected action defined by a previously recorded
expectation, the framework records E\textsuperscript{+} when supporting evidence
is explicitly recorded, or derives E\textsuperscript{-} only when absence is
verified under the conditions defined by Janus Core Lite v0.2.

\medskip
\noindent\textbf{Invariant 2 (Omission Detectability):}
An omission is detectable only relative to a previously recorded, versioned
expectation, a deterministic scope and cutoff, the applicable schema context,
and a reproducible observation.

\medskip
\noindent\textbf{Invariant 3 (Human Authority Closure):}
When an outcome requires explicit human accountability, the governance state is
not complete until a \texttt{HUMAN\_DECISION} is recorded. A recorded
\texttt{OMISSION\_DETECTED} event does not substitute for that decision.

\medskip
\noindent\textbf{Invariant 4 (Deterministic Reconstruction):}
Given the same log records, relevant schemas, and referenced evidence artifacts
up to the same cutoff, this framework profile requires reconstruction to produce
the same governance-relevant views and findings. This strengthens the Core's
rebuildability recommendation but does not define cross-log ordering.

\section{Runtime Observability}

Janus emits structured events at runtime, each containing:
\begin{itemize}
    \item Event type
    \item Timestamp
    \item Payload
    \item Error traces
\end{itemize}

This structured telemetry can support failure detection and end-to-end
traceability. Runtime telemetry is implementation evidence; it is not itself a
canonical governance event unless it satisfies the Core event contract.

\section{Evidence-Based Governance Benchmark}

\subsection{Metric Definitions}

\noindent\textbf{Event Capture Rate (ECR):}
\begin{equation}
  ECR = \frac{\text{Operations with required evidence}}{\text{Eligible governed operations}}
\end{equation}
ECR is evaluated within a declared scope and cutoff. The denominator contains
only operations to which a recorded protocol expectation applies; the numerator
contains those operations for which the evidence required by that expectation
is resolvably recorded. A value of 1.0 establishes coverage only within that
declared scope, not conformance of the system as a whole.

\medskip
\noindent\textbf{Governance Recording Interval (GRI):}
\begin{equation}
\begin{split}
  GRI = \; & t_{recorded}(\mathit{HUMAN\_DECISION}) \\
            & - \; t_{recorded}(\mathit{OMISSION\_DETECTED})
\end{split}
\end{equation}
GRI applies only when a \texttt{HUMAN\_DECISION} resolvably identifies the
preceding \texttt{OMISSION\_DETECTED} as an object. It measures the interval
between incorporation of two records. It measures neither omission-detection
latency nor the duration of human deliberation, and no Core rule requires such a
decision for every omission.

\medskip
\noindent\textbf{Protocol Violation Detectability Rate (PVDR):}
\begin{equation}
  PVDR = \frac{\text{Detected Violations}}{\text{Injected Violations}}
\end{equation}
The numerator and denominator are integer counts from a versioned injection
manifest. Each trial must identify the injected condition, expected detector,
scope, cutoff, observation artifact, and result. Aggregate values are not
reportable without the trial-level manifest.

\subsection{Evaluation Principle}

\begin{equation}
  \mathit{Governance\_Conformance} =
    \frac{\text{Observed Events}}{\text{Expected Protocol Events}}
\end{equation}

Expected protocol events are derived from versioned protocol expectations
recorded before observation. The observation scope identifies sources, subject,
filters, window, cutoff, procedure, and applicable schemas.

\section{Status of Empirical Validation}

No raw dataset or execution record supporting the numerical results reported in
v1.0 was located during reconciliation. Those values are withdrawn from this
revision. In particular, a PVDR of 0.95 cannot result from exactly 50 binary
trials. Future empirical results must publish the versioned protocol, injection
manifest, raw observations, calculation procedure, and execution environment.
Until those artifacts are available, the benchmark in this paper is a proposed
measurement specification rather than a validated or replicated result.

\section{Related Work}

Janus relates to several established paradigms and frameworks:

\begin{itemize}
    \item Event Sourcing~\cite{fowler}
    \item Command Query Responsibility Segregation (CQRS)~\cite{young}
    \item Distributed event ordering~\cite{lamport}
    \item Observability frameworks~\cite{otel}
    \item AI risk management~\cite{nist}
    \item Requirements traceability~\cite{gotel}
\end{itemize}

Janus differs from these approaches by integrating governance, evidence, and
human authority as first-class runtime constructs, rather than as post-hoc
auditing mechanisms.

\section{Limitations}

Current limitations of the Janus model include:

\begin{itemize}
    \item Lack of automated enforcement of authority boundaries
    \item Partial observability in external or third-party systems
    \item Human authority scalability challenges at high decision volumes
    \item No preserved raw dataset for the v1.0 numerical claims
\end{itemize}

The system currently focuses on detection and signaling of governance deviations,
rather than enforcing them automatically.

\subsection{Threats to Validity}
\begin{itemize}
    \item No independent replication by external teams
    \item Benchmark definitions require versioned scopes and trial manifests
    \item Controlled fault injection may not reflect production conditions
\end{itemize}

\section{Conclusion}

This paper proposes that governance can be evaluated through measurable
conformance between versioned expectations and observed behavior. Treating
evidence capture, omission detection, and human authority as first-class
operational constructs provides a testable foundation. Demonstrating
reproducibility requires publication and independent evaluation of the evidence
artifacts specified above.


\begin{thebibliography}{7}

\bibitem{foundational}
M. N. Sanchez Morales, ``Janus: A Minimal Governance Kernel for Deterministically Reconstructible Human--AI Development Systems,'' 2026. [Online]. Available: \url{https://doi.org/10.5281/zenodo.18974356}


\bibitem{fowler}
M. Fowler, ``Event Sourcing,'' 2005. [Online].
Available: \url{https://martinfowler.com/eaaDev/EventSourcing.html}


\bibitem{young}
G. Young, ``CQRS Documents,'' 2010. [Online].
Available: \url{https://cqrs.files.wordpress.com/2010/11/cqrs_documents.pdf}


\bibitem{lamport}
L. Lamport, ``Time, Clocks, and the Ordering of Events in a Distributed System,''
\emph{Commun. ACM}, vol.~21, no.~7, pp.~558--565, Jul. 1978.


\bibitem{otel}
OpenTelemetry, ``OpenTelemetry Specification,''
Cloud Native Computing Foundation (CNCF), 2023. [Online].
Available: \url{https://opentelemetry.io/docs/specs/otel/}


\bibitem{nist}
NIST, ``AI Risk Management Framework (AI RMF 1.0),'' 2023. [Online].
Available: \url{https://www.nist.gov/itl/ai-risk-management-framework}


\bibitem{gotel}
O. Gotel and A. Finkelstein, ``An Analysis of the Requirements Traceability
Problem,'' \emph{Proc. IEEE Int. Conf. Requirements Engineering}, 1994.

\end{thebibliography}

\end{document}
